\subsubsection{Итерация [[K]]}

\begin{enumerate}
  \item Найдём значение градиента $\omega^{([[K]])}=\textrm{grad}\,f(x^{([[K_PREV]])})$:
  \[
  \begin{gathered}
    \omega^{([[K]])}=(2ax_1^{([[K_PREV]])}+cx_2^{([[K_PREV]])}+d,2bx_2^{([[K_PREV]])}+cx_1^{([[K_PREV]])}+e)^T=\\
    =(2\cdot [[A]]\cdot [[X_PREV]]+[[C]]\cdot [[Y_PREV]]+[[D]], 2\cdot [[B]]\cdot [[Y_PREV]]+[[C]]\cdot [[X_PREV]]+[[E]])^T=[[OMEGA]]
  \end{gathered}
  \]

  \item Нормируем антиградиент \( \omega^{([[K]])} \), чтобы получить орт:
  \[
    u^{([[K]])}=-\dfrac{\omega^{([[K]])}}{|\omega^{([[K]])}|}=-\dfrac{[[OMEGA]]}{[[OMEGA_NORM]]}=([[U_X]],[[U_Y]])
  \]

  \item Найдём \( \beta^{([[K]])} \) как вершину параболы:
  \[
  \begin{gathered}
    \beta^{([[K]])}_1=\textrm{arg}\,\min_{\beta}f(x^{([[K_PREV]])}+\beta u^{([[K]])})=\\
    =-\dfrac{2ax^{([[K_PREV]])}_1u_1^{([[K]])}+2bx_2^{([[K_PREV]])}u_2^{([[K]])}+c(x_1^{([[K]])}u_2^{([[K]])}+x^{([[K_PREV]])}_2u_1^{([[K]])})+du_1^{([[K]])}+eu_2^{([[K]])}}{2(a(u_1^{([[K]])})^2)+b(u_2^{([[K]])})^2+cu_1^{([[K]])}u_2^{([[K]])}}=\\
    =-\dfrac{2\cdot [[A]]\cdot [[X_PREV]]\cdot [[U_X]]+\dots+[[D]]\cdot [[U_X]]+[[E]]\cdot [[U_Y]]}{2\cdot ([[A]]\cdot([[U_X]])^2+[[B]]\cdot ([[U_Y]])^2+[[C]]\cdot [[U_X]]\cdot [[U_Y]])}=[[BETA]]
  \end{gathered}
  \]

  \item Вычислим \( x^{([[K]])} \):
  \[
    x^{([[K]])}=x^{([[K_PREV]])}+\beta^{([[K]])}\cdot u^{([[K]])}=([[X_NEW]],[[Y_NEW]])^T
  \]

  \item Графическая интерпретация:
  \begin{figure}[H]
  \centering
  \begin{tikzpicture}
    \begin{axis}[
      width=10cm, height=10cm,
      axis lines=middle,
      xlabel={$x$},
      ylabel={$y$},
      xlabel style={right},
      ylabel style={above},
      colormap={blackwhite}{gray(0cm)=(0.4); gray(1cm)=(0.9)},
      view={0}{90}
    ]
      \addplot3[
        contour lua={
          levels={[[Z_LEVELS]]},
          labels=false,
        },
        samples=100,
        domain=[[X_MIN]]:[[X_MAX]],
        domain y=[[Y_MIN]]:[[Y_MAX]],
      ] {[[A]]*x^2 + [[B]]*y^2 + [[C]]*x*y + [[D]]*x + [[E]]*y + [[F]]};

      [[STEEPEST_PATH_ITER]]
        
    \end{axis}
  \end{tikzpicture}
  \caption{Итерация №[[K]] метода наискорейшего спуска.}
  \end{figure}

  \item Проверка критерия останова:
  \[
    [[BREAKPOINT_CHECK]]
  \]
\end{enumerate}
